\section{Introduction}
\label{sec:introduction}

% Follows the plan of the first page in outline.md; the first page must state how the work is relevant to real-world
% privacy. The last three paragraphs were written after the main results (24 October 2026).

Keystroke logs record how a text was written: every press and release, the pauses between them, the revisions and the
movements of the cursor. Writing research uses them to study the process of writing, and public corpora release logs
together with their texts and scores~\cite{tian2025klicke}. A~log also carries the text it produced, so a release
that is meant to keep the text apart from the process has to hide the text. One public release did this by masking
letters and digits:
the logs of a Kaggle competition on writing quality replaced every letter and digit by \texttt{q}, with the stated goal
``to prevent reproduction of the essay text''~\cite{kaggle2023lwp}. The lengths of words, the punctuation, the times
and the cursor positions remained, and the winning solution rebuilt the masked essays from the logs and used their
word and sentence lengths to predict their scores~\cite{inubushi2024first}.

The texts of such essays do not disappear when the logs are released. They were collected and scored, so they are held by
those who collected and scored them, and texts on the same prompts can be public; in a classroom, a teacher holds the
texts of a class. The authors of a public corpus, KLiCKe~\cite{tian2025klicke}, told us that the competition data were
drawn from that corpus before it was released with its texts; we did not check whether the released corpus contains the hidden essays
(Section~\ref{sec:threat-model:settings}). A party
who holds candidate texts does not need a labelled sample of a writer's keystrokes, which earlier work on releasing
keystroke data assumed an attacker would need~\cite{edwards2023csedm}; it can compare what the processed log leaves with
each text. This paper asks how far processing keeps logs apart from their texts against such a party, which is how we read the
stated goal, and what each form of processing costs research.

We answer this for one processed log against $N$ candidate texts on the same prompt, in a closed world, with $N =
10$, $100$ and $1{,}000$, on a public corpus whose texts are known (KLiCKe). The forms are the masking used in the
competition and alternatives that release less: key categories, merged separators, rounded times, segment values
alone and random suppression of presses. The attackers know the processing and range from one that uses only the
total length to one that rebuilds the final draft and one that scores the press intervals. For each form we report
the strongest attacker's accuracy, as a multiple of $1/N$ and as a gain over the total length, and how far ten
standard process measures of writing research can still be computed; the same kind of matching is also run on
Japanese translation logs typed with an input method. The study was preregistered before any data were acquired.

The main harm we consider is the link between a processed log and its text for a party that holds the text: the times,
the revisions and any released score of the log are then linked to that text. Where the writer of a text is known, as
in a class or for an essay its writer has published, the log would also be linked to a named person; we know of no
release in which this has happened and state it only as a consequence (Section~\ref{sec:threat-model}).

Against the attackers tried, the masking kept no log apart from its text: an attacker that rebuilds the final draft,
as the winning solution of the competition did, chose the true essay among 1,000 for every log (C1,
Section~\ref{sec:results:masking}). Forms that release less lowered matching only at a cost to the process measures:
keeping the key categories still left 86.7\% at $N = 1{,}000$, every form that kept the measures close to their
unprocessed values (a median correlation of 0.84 or more) left 70\% or more, and the forms that brought matching near the total length kept them at a
median correlation of at most 0.70 or not at all (C2, Sections~\ref{sec:results:forms} and~\ref{sec:results:utility}). We frame a threat in which the
auxiliary information is the text, held outside the release, and give the settings in which it arises (C3,
Section~\ref{sec:threat-model}). The results describe the attackers tried; they are not bounds on the risk.
